\section{Introducción}\label{sec:introduccion}

Una parte de la energía que entregan las empresas distribuidoras se consume sin pagarse.
Son las pérdidas no técnicas: conexiones directas a la red, medidores manipulados y lecturas
o registros alterados \parencite{glauner2017}. A diferencia de las pérdidas técnicas, que se
disipan en el transporte, estas tienen un origen identificable en el cliente, y su costo
termina recayendo sobre los usuarios que sí pagan.

En Colombia el problema se concentra en la región Caribe, que en junio de 2024 representó
el 28,5\,\% de la demanda nacional de energía eléctrica, según datos de XM
\parencite{elheraldo2024}. En los barrios eléctricamente subnormales que atiende Afinia, la
empresa recauda apenas el 5\,\% de lo que factura, lo que equivale a pérdidas superiores a
300.000 millones de pesos al año \parencite{portafolio2025}. En Air-e, el recaudo acumulado
de 2025 fue del 74,55\,\%, frente a un promedio nacional del 95\,\% \parencite{elheraldo2025}.
Las cifras de recaudo mezclan el fraude con la simple falta de pago, pero dan una idea de la
escala del problema en las distribuidoras de la región.

La respuesta habitual de las empresas es la inspección en campo. Cada visita tiene un costo
y la capacidad de inspección es finita, así que la pregunta operativa es a quién visitar
primero. Por eso el problema se plantea aquí como una tarea de
clasificación orientada a la priorización: ordenar a los clientes según su probabilidad de
fraude a partir de la única información que la empresa tiene de todos ellos, el historial
de facturación y los atributos del contrato.

Ese planteamiento choca con una dificultad que la literatura suele pasar por alto. La
etiqueta de fraude solo existe para los clientes que fueron inspeccionados y en los que se
confirmó el fraude, de modo que mide fraude \emph{detectado}. Un cliente con un historial
largo ha tenido más ocasiones de ser inspeccionado, y una inspección exitosa deja rastros en
el propio historial, como el reemplazo del contador. Un modelo puede aprender esos rastros y
obtener cifras altas en validación sin capacidad real de anticipar el fraude en clientes que
todavía no han sido visitados, que es justamente el uso para el que se construye.

El estudio usa los datos públicos de la Société Tunisienne de l'Électricité et du Gaz
(STEG), liberados en una competencia de la plataforma Zindi \parencite{zindi2019}. Es uno de
los pocos conjuntos abiertos con etiqueta de fraude por cliente y con el historial de
facturación de contadores convencionales, la misma información disponible para las
distribuidoras que no cuentan con medición inteligente. El aporte que se busca tiene dos
partes: un modelo que priorice inspecciones con una mejora medible sobre el azar y sobre una
regla de exposición, y una evaluación que separe la señal de comportamiento del cliente de
la huella que deja la inspección.
